% docs/manual/operations/appendices/C-cache-layout.tex
% 附录 C：缓存目录布局

\chapter{缓存目录布局}

\openbench{} 输出目录 \file{output/<run>/} 完整结构。

\section{完整目录树}

\begin{minted}{text}
output/<run>/
├── config.yaml                  # 本次评估的解析后完整配置
├── run.log                      # 运行日志（INFO 与上）
├── data/                        # zarr 缓存（中间数据，可重建）
│   ├── <task_key>.zarr/         # 每 task 一个 zarr group
│   │   ├── sim/
│   │   ├── ref/
│   │   ├── .zgroup
│   │   ├── .zarray
│   │   └── ...
│   └── ...
├── metrics/                     # 评估指标 CSV
│   ├── <var>/
│   │   ├── <sim>_metrics.csv
│   │   └── ...
│   └── ...
├── scores/                      # 归一化打分
│   └── <var>/<sim>_scores.csv
├── figures/                     # 全部可视化（按 var × kind 子目录）
│   └── <var>/<kind>/<sim>_<ref>.png
├── comparisons/                 # 多 sim 对比（仅 comparison.enabled=true）
│   ├── portrait_table.csv
│   ├── portrait_plot.png
│   ├── radar_plot.png
│   ├── ranking.csv
│   └── ranking_table.png
├── statistics/                  # 仅 statistics.enabled=true
│   ├── ANOVA/
│   ├── Correlation/
│   └── ...
├── reports/                     # HTML/PDF 总结报告
│   ├── report.html
│   ├── report.pdf
│   └── sections/                # 子报告（按变量分）
├── scratch/                     # 任务级 cache + 临时
│   ├── cache/
│   │   ├── <task_key>.json
│   │   └── ...
│   └── ...
├── tmp/                         # 单次运行临时（每次清）
└── debug/                       # 仅 --dump-config
    ├── runner_config.yaml
    ├── main_nl.yaml
    ├── ref_nml.yaml
    ├── sim_nml.yaml
    └── fig_nml.yaml
\end{minted}

\section{每层用途}

\begin{longtable}{p{2.5cm} p{4cm} p{6cm}}
  \toprule
  目录 & 用途 & 可删？ \\
  \midrule
  \endhead
  \texttt{config.yaml} & 解析后的完整配置（含默认） & 否（追溯用） \\
  \texttt{run.log} & 完整日志 & 旧的可压缩归档 \\
  \texttt{data/} & zarr 中间产物 & 可（重建费时） \\
  \texttt{metrics/} & 指标 CSV & \emph{结论数据，不要删} \\
  \texttt{scores/} & 打分 CSV & 同上 \\
  \texttt{figures/} & PNG 图 & 可（重出图较快） \\
  \texttt{comparisons/} & 对比表与图 & \emph{结论；最好保留} \\
  \texttt{statistics/} & 统计输出 & \emph{结论；最好保留} \\
  \texttt{reports/} & HTML/PDF & 可（从 metric/figure 重建） \\
  \texttt{scratch/cache/} & task 级 hash & 可（让下次重算） \\
  \texttt{tmp/} & 临时 & 可（每次 run 自动清） \\
  \texttt{debug/} & 中间 namelist & 可 \\
  \bottomrule
\end{longtable}

\section{磁盘占用经验值}

以全球 monthly × 5 年 × 1 sim × 4 vars 评估为例：

\begin{longtable}{p{3cm} p{3cm} p{6cm}}
  \toprule
  目录 & 大小 & 备注 \\
  \midrule
  \endhead
  data/ & 200 MB - 2 GB & zarr 压缩；与变量数线性 \\
  metrics/ & 几 MB & CSV 小 \\
  scores/ & 几 MB & 同上 \\
  figures/ & 50-200 MB & 按变量与图种数 \\
  comparisons/ & 10-50 MB & 多 sim 时增加 \\
  statistics/ & 50-300 MB & 统计模块多时增 \\
  reports/ & 50-300 MB & 嵌入图 \\
  scratch/cache/ & < 1 MB & JSON \\
  \bottomrule
\end{longtable}

总：单次 monthly 全球评估约 500 MB - 3 GB。daily/hourly + 多 sim 可达数十 GB。

\section{升级时哪些保留}

\openbench{} 升级后：

\begin{itemize}
  \item \emph{保留}：metrics、scores、comparisons、statistics、figures、reports（结论数据）
  \item \emph{清理}：data、scratch（cache 可能与新版不兼容）
  \item \emph{重生成}：reports/figures（如果想用新版样式）
\end{itemize}

\section{多用户共享布局示例}

\begin{minted}{text}
/scratch/openbench-shared/
├── refs/                # 共享只读 reference 数据集
│   ├── Grid/
│   └── Station/
├── eval-2026-Q1/
│   ├── alice-eval1/     # alice 的 evaluation
│   ├── alice-eval2/
│   ├── bob-eval1/
│   └── ...
└── archived/            # 旧 evaluation 归档
    └── 2025-Q4/
\end{minted}

每用户 evaluation 独立目录；共享数据走 \texttt{refs/}。
